\subsection{محدودیت‌ها و تحلیل شکست}
تحلیل ریشه‌ای شکست مدل در بنچ‌مارک \lr{RoBERTa} به عوامل زیر بازمی‌گردد:
\begin{itemize}
    \item \textbf{ناسازگاری ماسک‌گذاری توجه در ساختار تو در توی \lr{RobertaSelfAttention}:} معماری \lr{RoBERTa} از ماسک‌های باینری با مقادیر بسیار منفی بزرگ (نظیر $-10000.0$) برای نشانه‌های پدینگ استفاده می‌کند. جایگزینی مستقیم تابع بیشینه‌ساز نرم در زیرماژول‌های تو در تو، در ترکیب با فرآیند استخراج برچسب طبقه‌بند بدون هماهنگی سر خروجی، منجر به عدم محاسبه صحیح برچسب در حلقه ارزیابی دقت و گزارش دقت \lr{0.00\%} با وجود محاسبه شاخص F1 شده است.
    \item \textbf{پیچیدگی تزریق لایه‌ای در معماری‌های با مکانیزم توجه ادغام‌شده:} در مدل‌هایی که تابع بیشینه‌ساز نرم به عنوان عملگر درونی در متد \lr{forward} پیاده‌سازی شده و فاقد ماژول مجزای \lr{nn.Softmax} هستند، فشرده‌ساز ناگزیر به جایگزینی کل بلوک توجه بوده که این امر در صورت کوچکترین عدم انطباق تانسورهای خروجی با خطای ساختاری همراه می‌شود.
\end{itemize}
